\documentclass[12pt]{article}
\usepackage[T1]{fontenc}
\usepackage[utf8]{inputenc}
\usepackage{lmodern}
\usepackage{textcomp}
\usepackage{enumitem}
\usepackage{geometry}
\usepackage{graphicx}
\usepackage{amsmath, amssymb}
\geometry{margin=1in}
\begin{document}
\begin{center}
Economics - Undergraduate Level Assessment 8 \
Total Marks: 40 \
Answer all questions.
\end{center}

\section*{Multiple Choice (10 marks)}
\begin{enumerate}[label=\arabic*.]
\item Explain why the opportunity for unions to increase the demand for labor is limited.
\begin{enumerate}[label=(\alph*)]
    \item Unions determine labor productivity
    \item Unions can directly control employment rates nationally
    \item Unions are the primary drivers of economic growth and thereby increase labor demand
    \item Unions can significantly increase labor demand
    \item Unions have limited means to influence the demand for labor
\end{enumerate}
\item The appropriate fiscal policy to remedy inflation calls for
\begin{enumerate}[label=(\alph*)]
    \item decreased government spending and increased taxes.
    \item the federal government to decrease the money supply.
    \item the federal government to run a deficit.
    \item decreased taxes and increased government spending.
    \item the federal government to increase the money supply.
\end{enumerate}
\item Which of the following examples would result in consumers paying for the largest burden of an excise tax placed on a producer?
\begin{enumerate}[label=(\alph*)]
    \item If both the demand and supply curves are price inelastic
    \item If the demand curve is perfectly elastic and the supply curve is perfectly inelastic
    \item If the demand curve is price inelastic and the supply curve is price elastic
    \item If both the demand and supply curves are price elastic
    \item If the demand curve is price elastic and the supply curve is price inelastic
\end{enumerate}
\item What is meant by the term marginal costs?
\begin{enumerate}[label=(\alph*)]
    \item the cost of all produced units
    \item the fixed cost associated with the production facility
    \item the cost of total production
    \item the cost of initial unit of production
    \item the additional cost of one more unit of production
\end{enumerate}
\item Which of the following is a characteristic of monopolistic competition?
\begin{enumerate}[label=(\alph*)]
    \item Very few competitors.
    \item One firm sets the market price.
    \item Identical products offered by each firm.
    \item Excess capacity.
    \item Significant barriers to entry.
\end{enumerate}
\end{enumerate}

\section*{True / False (10 marks)}
\textit{Answer True or False}
\begin{enumerate}[label=\arabic*.]
\setcounter{enumi}{5}
\item True or False: The correct answer to 'Which one of the following statements best describes the algebraic representation of the fitted regression line?' is 'y\_t = \ensuremath{\hat{\alpha}} + \ensuremath{\hat{\beta}}x\_t'.
\item True or False: A monopoly's marginal revenue curve is the same as its demand curve.
\item True or False: The fixed effects panel model is commonly referred to as the random effects model in statistical literature.
\item True or False: An increase in printing costs will shift the supply curve for textbooks to the left, indicating decreased supply.
\item True or False: A nation producing inside its production possibility curve is likely experiencing an economic recession.
\end{enumerate}

\section*{Long Form Response (20 marks)}
\textit{Answer each question with a clear and complete explanation.}
\begin{enumerate}[label=\arabic*.]
\setcounter{enumi}{10}
\item Discuss the factors that influence the decision to hire labor in a market where the wage is set at \$4.50, and analyze the implications of hiring 8 units of labor.
\item Discuss the disadvantages of using pure time-series models in economic analysis, particularly in terms of their theoretical motivation compared to structural models.
\end{enumerate}

\vfill
\noindent\textit{End of Paper}
\end{document}